\section{Experiment 3}
\label{Exp3}
\setlength{\parindent}{20pt}
\justifying

\noindent
Presenting the characters in the BACS-1 font may have expanded the horizon of participant recruitment, but the unfamiliarity of the characters may have required additional effort to process and learn. This took away cognitive resources from the learning of higher-level statistics, and prevented observable co-occurrence learning to occur. Ideally, the authors' aim would be to develop a version of this paradigm with mean 2M accuracy above chance in order to see more meaningful variability in scores. One further drawback from the stimuli in previous versions of the experiment was the choice to create stimuli for the two languages which were only partially overlapping in the characters used to form them. This was selected to make stimuli ecological and similar to European languages, for example: most of these use the standard Latin alphabet, with some language-specific additions (``é'', ``ô'', ``\ss'', ``œ''). However, this gives participants some indication as to which language morphemes belong to, beyond pure co-occurrence statistics as intended. 

For these reasons, the experiment was adjusted to present stimuli in Latin characters fully overlapping across both languages.

\subsection{Methods}
\label{Methods3}

\subsubsection{\large{\emph{Participants}}}
\label{Participants3}

\noindent
In order to present stimuli in Latin characters while retaining the novelty of the strings presented to participants, the range of languages spoken by participants had to be limited. A set of languages were chosen to keep the variety of language backgrounds as wide as possible. However, researchers had to ensure that a string presented in the experiment was not a real lexical item in any of the languages spoken by participants. To this end, the chosen languages had to be similar to each other to prevent the exclusion of too many random strings that might be used in the experiment. These languages were selected as English, German, Dutch, Spanish, Portuguese, French, Catalan, and Italian.

Therefore, the participant sample size and criteria remained identical to the previous two experiments, with the new restriction of them speaking only the languages previously mentioned, or any other language not using a Latin script.


\subsubsection{\large{\emph{Stimuli}}}
\label{Stimuli3}

\noindent
Given the necessity to generate morpheme strings that are not real lexical items in any of the above listed languages, a pipeline was developed for their generation. All possible 3- to 5-character combinations of all the letter of the Latin alphabet were assembled, and any item present in these languages was excluded. From the remaining sample, 600 were randomly selected, and replaced by new random choices if any are at a Levenshtein distance of less than 2 or have 2-character chunks within them in common with another string. This leaves a pool of 572 strings, from which each participant gets a random selection of the 30 morphemes to be used in their submission. The result of this new stimulus generation procedure was that the morphemes for both language 1 and language 2 contain all (unaccented) letters of the Latin alphabet, and are therefore completely indistinguishable at the letter level. 


\subsubsection{\large{\emph{Procedure}}}
\label{Procedure3}

\noindent
Given that Latin characters are easier to process than BACS-1 characters, this may lead to an over-familiarity with the presented morphemes, thereby causing participants to select all 2M combinations as congruent. A large ``yes'' bias would greatly diminish participants' ability to reject between-language combinations. Indeed, a pilot experiment with 10 participants revealed that the c value of participants' responses to the 2M condition was -0.72. To correct this over-familiarity, the morpheme familiarity test might be removed. A pilot was done under these conditions, and yielded a much more acceptable c of -0.34, reassuringly in line with previous experiments in this paper (experiment 1 c = -0.37; experiment 1 c = -0.35). The mean familiarity score for this pilot was also 57\%, in line with both previous versions. This choice was therefore chosen as giving researchers the closest match to previous versions and their ``yes'' biases, and offered the best chance for variability in responses. Experiment 3's procedure was therefore identical to experiment 1.

\subsubsection{\large{\emph{Analysis}}}
\label{Analysis3}

\noindent
Data analysis was identical to experiment 2. 